\begin{explanationbox}
\textbf{\textcolor{CExplanation}{一句话直觉}}
把每个变量的平方“分摊”给下一个变量作分母，这样的循环分式之和，不会小于变量自身之和。它描述了一种“循环放大”效应，但总和仍有下限。

\medskip
\textbf{\textcolor{CExplanation}{核心拆解}}
\begin{description}[style=nextline, leftmargin=2.8em, labelsep=0.5em, itemsep=0.45em]
    \item[\textbf{要点一（结构对称）：}] 不等式左边三个分式呈 $x \to y \to z \to x$ 的循环对称结构。
    \item[\textbf{要点二（分式降幂）：}] 左边每一项是二次除以一次，整体是齐一次式，右边也是齐一次式。
\end{description}

\medskip
\textbf{\textcolor{CExplanation}{几何本质（旋转放缩）}}
可以将 $\frac{x^2}{y}$ 视为以 $x$ 和 $\frac{x}{y}$ 为两边的矩形面积的一种度量。不等式表明，经过循环“旋转”和“伸缩”后的面积之和，总不小于原始边长之和。当三个变量相等时，这种“旋转伸缩”效应消失，达到平衡。

\medskip
\textbf{\textcolor{CExplanation}{代数意义（柯西变形）}}
从代数角度看，该不等式是柯西-施瓦茨不等式一个精妙的变形。通过将 $x+y+z$ 拆分为 $\sqrt{x}\sqrt{\frac{x}{y}} \cdot \sqrt{y} + \dots$ 等形式，并应用 $a^2+b^2+c^2 \geq ab+bc+ca$，可自然导出结论。它揭示了平方和与交叉项和之间的深层关系。

\medskip
\textbf{\textcolor{CExplanation}{考点价值（综合应用）}}
\begin{itemize}[leftmargin=*, itemsep=0.5em, topsep=0.4em]
    \item \textbf{考法一（直接证明求最值）：} 作为已知结论直接用于证明其他不等式或求多元表达式的最小值。
    \item \textbf{考法二（变形识别应用）：} 题目可能将 $\frac{x^2}{y}$ 变形为 $\frac{x^3}{xy}$ 等形式，需要逆向识别出本结构方可应用。
\end{itemize}

\medskip
\textbf{\textcolor{CExplanation}{顿悟点}}
\textbf{关键在于发现 $\frac{x^2}{y} + y \geq 2x$，并对三个这样的不等式进行循环求和。} 这是将非对称分式化为对称整体最常用的技巧。

\medskip
\textbf{\textcolor{CExplanation}{使用场景}}
\begin{itemize}[leftmargin=*, itemsep=0.5em, topsep=0.4em]
    \item \textbf{场景一（循环分式求和）：} 当遇到 $\frac{a^2}{b}+\frac{b^2}{c}+\frac{c^2}{a}$ 这类循环分式求和求最小值时，优先考虑此结论。
    \item \textbf{场景二（条件最值证明）：} 在已知 $x,y,z>0$ 且满足某条件（如 $xyz=1$）时，证明包含此类循环分式的不等式。
\end{itemize}
\end{explanationbox}
